\documentclass[10pt,a4paper]{amsart}
\usepackage{amsmath,amssymb,amsthm,graphicx}
\usepackage{verbatim}
\newenvironment{reference}{\comment}{\endcomment}
%\newenvironment{reference}{}{}
\newenvironment{slogan}{\comment}{\endcomment}
\newenvironment{history}{\comment}{\endcomment}

% For commutative diagrams we use Xy-pic
\usepackage[all]{xy}

% We use 2cell for 2-commutative diagrams.
\xyoption{2cell}
% Theorem environments.
%
\theoremstyle{plain}
\newtheorem{theorem}[subsection]{Theorem}
\newtheorem{proposition}[subsection]{Proposition}
\newtheorem{lemma}[subsection]{Lemma}

\theoremstyle{definition}
\newtheorem{definition}[subsection]{Definition}
\newtheorem{example}[subsection]{Example}
\newtheorem{exercise}[subsection]{Exercise}
\newtheorem{situation}[subsection]{Situation}

\theoremstyle{remark}
\newtheorem{remark}[subsection]{Remark}
\newtheorem{remarks}[subsection]{Remarks}

\numberwithin{equation}{subsection}

% Macros
%
\def\lim{\mathop{\mathrm{lim}}\nolimits}
\def\colim{\mathop{\mathrm{colim}}\nolimits}
\def\Spec{\mathop{\mathrm{Spec}}}
\def\Hom{\mathop{\mathrm{Hom}}\nolimits}
\def\Ext{\mathop{\mathrm{Ext}}\nolimits}
\def\SheafHom{\mathop{\mathcal{H}\!\mathit{om}}\nolimits}
\def\SheafExt{\mathop{\mathcal{E}\!\mathit{xt}}\nolimits}
\def\Sch{\mathit{Sch}}
\def\Mor{\mathop{\mathrm{Mor}}\nolimits}
\def\Ob{\mathop{\mathrm{Ob}}\nolimits}
\def\Sh{\mathop{\mathit{Sh}}\nolimits}
\def\NL{\mathop{N\!L}\nolimits}
\def\CH{\mathop{\mathrm{CH}}\nolimits}
\def\proetale{{pro\text{-}\acute{e}tale}}
\def\etale{{\acute{e}tale}}
\def\QCoh{\mathit{QCoh}}
\def\Ker{\mathop{\mathrm{Ker}}}
\def\Im{\mathop{\mathrm{Im}}}
\def\Coker{\mathop{\mathrm{Coker}}}
\def\Coim{\mathop{\mathrm{Coim}}}

% Boxtimes
%
\DeclareMathSymbol{\boxtimes}{\mathbin}{AMSa}{"02}

%
% Macros for moduli stacks/spaces
%
\def\QCohstack{\mathcal{QC}\!\mathit{oh}}
\def\Cohstack{\mathcal{C}\!\mathit{oh}}
\def\Spacesstack{\mathcal{S}\!\mathit{paces}}
\def\Quotfunctor{\mathrm{Quot}}
\def\Hilbfunctor{\mathrm{Hilb}}
\def\Curvesstack{\mathcal{C}\!\mathit{urves}}
\def\Polarizedstack{\mathcal{P}\!\mathit{olarized}}
\def\Complexesstack{\mathcal{C}\!\mathit{omplexes}}
% \Pic is the operator that assigns to X its picard group, usage \Pic(X)
% \Picardstack_{X/B} denotes the Picard stack of X over B
% \Picardfunctor_{X/B} denotes the Picard functor of X over B
\def\Pic{\mathop{\mathrm{Pic}}\nolimits}
\def\Picardstack{\mathcal{P}\!\mathit{ic}}
\def\Picardfunctor{\mathrm{Pic}}
\def\Deformationcategory{\mathcal{D}\!\mathit{ef}}
\usepackage{cleveref}
\crefname{originalreference}{original reference}{original references}
\Crefname{originalreference}{Original reference}{Original references}
\makeatletter
\newlabel{lemma-exact-sequence}{{0AXF}{source}{}{sourcepointer-1}{}}
\newlabel{lemma-exact-sequence@cref}{{[originalreference][0][]0AXF}{[1][0][]source}{}{}{}}
\makeatother
\crefname{originallemma}{original lemma}{original lemmas}
\Crefname{originallemma}{Original lemma}{Original lemmas}
\makeatletter
\newlabel{limits-lemma-modifications}{{0B3X}{source}{}{sourcepointer-2}{}}
\newlabel{limits-lemma-modifications@cref}{{[originallemma][0][]0B3X}{[1][0][]source}{}{}{}}
\makeatother
\crefname{originallemma}{original lemma}{original lemmas}
\Crefname{originallemma}{Original lemma}{Original lemmas}
\makeatletter
\newlabel{limits-lemma-modifications-properties}{{0BFN}{source}{}{sourcepointer-3}{}}
\newlabel{limits-lemma-modifications-properties@cref}{{[originallemma][0][]0BFN}{[1][0][]source}{}{}{}}
\makeatother
\crefname{originallemma}{original lemma}{original lemmas}
\Crefname{originallemma}{Original lemma}{Original lemmas}
\makeatletter
\newlabel{more-morphisms-lemma-category-projective}{{0B46}{source}{}{sourcepointer-4}{}}
\newlabel{more-morphisms-lemma-category-projective@cref}{{[originallemma][0][]0B46}{[1][0][]source}{}{}{}}
\makeatother
\crefname{originallemma}{original lemma}{original lemmas}
\Crefname{originallemma}{Original lemma}{Original lemmas}
\makeatletter
\newlabel{morphisms-lemma-algebraic-residue-field-extension-closed-point-fibre}{{01TE}{source}{}{sourcepointer-5}{}}
\newlabel{morphisms-lemma-algebraic-residue-field-extension-closed-point-fibre@cref}{{[originallemma][0][]01TE}{[1][0][]source}{}{}{}}
\makeatother
\crefname{originallemma}{original lemma}{original lemmas}
\Crefname{originallemma}{Original lemma}{Original lemmas}
\makeatletter
\newlabel{morphisms-lemma-dimension-formula}{{02JU}{source}{}{sourcepointer-6}{}}
\newlabel{morphisms-lemma-dimension-formula@cref}{{[originallemma][0][]02JU}{[1][0][]source}{}{}{}}
\makeatother
\crefname{originallemma}{original lemma}{original lemmas}
\Crefname{originallemma}{Original lemma}{Original lemmas}
\makeatletter
\newlabel{morphisms-lemma-nagata-normalization}{{035S}{source}{}{sourcepointer-7}{}}
\newlabel{morphisms-lemma-nagata-normalization@cref}{{[originallemma][0][]035S}{[1][0][]source}{}{}{}}
\makeatother
\title{Rationality of a normal surface singularity propagates along essentially finite local extensions}
\author{The Stacks Project Authors}
\date{stacks-project:resolve.tex:local-snapshot-20261012; GNU FDL 1.2 or later}
\begin{document}
\maketitle
\noindent\textit{Editorial note.} All delivered mathematical text is the author's own, copied from the pinned Stacks Project chapter file resolve.tex; no statement or proof step was rewritten, reordered or removed.\par\medskip
\noindent\textit{Editorial note.} The counted proof is the author's complete argument: the realization of the extension as a local ring of a projective modification at a closed point, the reduction to a normal projective modification with the same local ring, and the spreading of modifications to the projective scheme with the transfer of the cohomological vanishing through the exact sequence.\par\medskip
\noindent\textit{Editorial note.} The label lemma-rational-propagates names the counted statement, and the proof environment immediately following it in the source is delivered verbatim as its proof.\par\medskip
\noindent\textit{Editorial note.} The counted statement is carried by the author's own lemma environment in the chapter source.\par\medskip
\noindent\textit{Editorial note.} Pre-open gate: python operations/agent-authoring/tools/precard\_gate.py --source resolve.tex --label lemma-rational-propagates --start 2034 --end 2084 --local-id 28076 returned RESULT=OK (L28076 free, occupancy FREE, public interval clear). Pre-assembly tag self-check: canonical tag resolved and not present in reference/public-items.jsonl.\par\medskip
\noindent\textit{Editorial note.} The statement and proof are copied from the pinned chapter file; the cited inputs are kept as external pointers with their canonical Stacks tags.\par\medskip
\noindent\textit{Allowed standard prerequisites (editorial).}\begin{itemize}
\item Realization of the extension as a local ring of a projective modification. Standard construction taking a finite type algebra over A presenting B and taking the closure of its spectrum in a projective space over A, together with the dimension formula and the algebraicity of residue field extensions for closed points. Supplies the projective modification X and the closed point whose local ring is B.
\item Finiteness of normalization for Nagata rings and projectivity of the normalized modification. Morphisms, the lemma that the normalization of a Nagata scheme is finite, and More on Morphisms, the lemma that the category of projective schemes is stable under composition. Supplies the reduction to a normal projective modification and the unique point above the chosen point with the same local ring.
\item Descent of modifications along a localization and vanishing of the first cohomology. Limits, the lemmas spreading a modification of the local ring to a modification of the projective scheme and checking properness, and Same chapter, the exact sequence used to conclude the vanishing of the first direct image. Supply the modification of the projective scheme restricting to the given one and the conclusion that its first cohomology vanishes on the modification of the local ring.
\end{itemize}\medskip
\ifdefined\hypertarget\hypertarget{sourcepointer-1}{}\fi
\noindent\textit{Source reference.} Original reference 0AXF: sources/stacks-resolve/pinned-original-source.tex, source lines 1901--1901. Prerequisite pointer invoked at proof line(s) 2082; it names a standard fact of the same chapter that the counted argument uses, and every step of the counted proof built on it is retained verbatim in the delivered body.\par\medskip
\ifdefined\hypertarget\hypertarget{sourcepointer-2}{}\fi
\noindent\textit{Source reference.} Original lemma 0B3X: sources/stacks-chow/tags.txt, source lines 13980--13980. Prerequisite pointer invoked at proof line(s) 2071; it names a standard fact of a sibling chapter that the counted argument uses, and every step of the counted proof built on it is retained verbatim in the delivered body.\par\medskip
\ifdefined\hypertarget\hypertarget{sourcepointer-3}{}\fi
\noindent\textit{Source reference.} Original lemma 0BFN: sources/stacks-chow/tags.txt, source lines 14391--14391. Prerequisite pointer invoked at proof line(s) 2075; it names a standard fact of a sibling chapter that the counted argument uses, and every step of the counted proof built on it is retained verbatim in the delivered body.\par\medskip
\ifdefined\hypertarget\hypertarget{sourcepointer-4}{}\fi
\noindent\textit{Source reference.} Original lemma 0B46: sources/stacks-chow/tags.txt, source lines 13989--13989. Prerequisite pointer invoked at proof line(s) 2062; it names a standard fact of a sibling chapter that the counted argument uses, and every step of the counted proof built on it is retained verbatim in the delivered body.\par\medskip
\ifdefined\hypertarget\hypertarget{sourcepointer-5}{}\fi
\noindent\textit{Source reference.} Original lemma 01TE: sources/stacks-chow/tags.txt, source lines 2371--2371. Prerequisite pointer invoked at proof line(s) 2055; it names a standard fact of a sibling chapter that the counted argument uses, and every step of the counted proof built on it is retained verbatim in the delivered body.\par\medskip
\ifdefined\hypertarget\hypertarget{sourcepointer-6}{}\fi
\noindent\textit{Source reference.} Original lemma 02JU: sources/stacks-chow/tags.txt, source lines 3349--3349. Prerequisite pointer invoked at proof line(s) 2052; it names a standard fact of a sibling chapter that the counted argument uses, and every step of the counted proof built on it is retained verbatim in the delivered body.\par\medskip
\ifdefined\hypertarget\hypertarget{sourcepointer-7}{}\fi
\noindent\textit{Source reference.} Original lemma 035S: sources/stacks-chow/tags.txt, source lines 4139--4139. Prerequisite pointer invoked at proof line(s) 2059; it names a standard fact of a sibling chapter that the counted argument uses, and every step of the counted proof built on it is retained verbatim in the delivered body.\par\medskip
\begin{lemma}
\label{lemma-rational-propagates}
Let $(A, \mathfrak m, \kappa)$ be a local normal Nagata domain of
dimension $2$ which defines a rational singularity. Let $A \subset B$
be a local extension of domains with the same fraction field
which is essentially of finite type such
that $\dim(B) = 2$ and $B$ normal. Then $B$ defines a rational singularity.
\end{lemma}

\begin{proof}
Choose a finite type $A$-algebra $C$ such that $B = C_\mathfrak q$
for some prime $\mathfrak q \subset C$. After replacing
$C$ by the image of $C$ in $B$ we may assume that $C$ is a domain
with fraction field equal to the fraction field of $A$.
Then we can choose a closed immersion $\Spec(C) \to \mathbf{A}^n_A$
and take the closure in $\mathbf{P}^n_A$ to conclude that $B$
is isomorphic to $\mathcal{O}_{X, x}$ for some closed point $x \in X$
of a projective modification $X \to \Spec(A)$.
(Morphisms, Lemma \ref{morphisms-lemma-dimension-formula},
shows that $\kappa(x)$ is finite over $\kappa$ and then
Morphisms, Lemma
\ref{morphisms-lemma-algebraic-residue-field-extension-closed-point-fibre}
shows that $x$ is a closed point.)
Let $\nu : X^\nu \to X$ be the normalization.
Since $A$ is Nagata the morphism $\nu$ is finite (Morphisms, Lemma
\ref{morphisms-lemma-nagata-normalization}).
Thus $X^\nu$ is projective over $A$ by
More on Morphisms, Lemma
\ref{more-morphisms-lemma-category-projective}.
Since $B = \mathcal{O}_{X, x}$ is normal, we see that
$\mathcal{O}_{X, x} = (\nu_*\mathcal{O}_{X^\nu})_x$.
Hence there is a unique point $x^\nu \in X^\nu$ lying over $x$
and $\mathcal{O}_{X^\nu, x^\nu} = \mathcal{O}_{X, x}$.
Thus we may assume $X$ is normal and projective over $A$.
Let $Y \to \Spec(\mathcal{O}_{X, x}) = \Spec(B)$
be a modification with $Y$ normal.
We have to show that $H^1(Y, \mathcal{O}_Y) = 0$. By
Limits, Lemma \ref{limits-lemma-modifications}
we can find a morphism of schemes $g : X' \to X$ which is an isomorphism
over $X \setminus \{x\}$ such that $X' \times_X \Spec(\mathcal{O}_{X, x})$
is isomorphic to $Y$. Then $g$ is a modification as it is proper by
Limits, Lemma \ref{limits-lemma-modifications-properties}.
The local ring of $X'$ at a point of $x'$ is either isomorphic
to the local ring of $X$ at $g(x')$ if $g(x') \not = x$ and
if $g(x') = x$, then the local ring of $X'$ at $x'$ is isomorphic
to the local ring of $Y$ at the corresponding point. Hence we see
that $X'$ is normal as both $X$ and $Y$ are normal.
Thus $H^1(X', \mathcal{O}_{X'}) = 0$ by our assumption on $A$.
By Lemma \ref{lemma-exact-sequence} we have $R^1g_*\mathcal{O}_{X'} = 0$.
Clearly this means that $H^1(Y, \mathcal{O}_Y) = 0$ as desired.
\end{proof}

\end{document}
